\chapter*{Abstract}
\addcontentsline{toc}{chapter}{Abstract}

This dissertation provides a comprehensive empirical investigation of the
\textit{Bitcoin Volatility Risk Premium} (BVRP), defined as the difference
between 30-day implied volatility (IV30D) and 30-day realized volatility
(RV30D). IV30D is sourced from the Deribit Volatility Index (DVOL), while
RV30D is constructed from daily log returns of Bitcoin spot prices. The sample
covers April 2021 to April 2026, yielding 1,776 daily observations. The study
examines the BVRP's statistical properties, its predictive content for future
returns, its out-of-sample predictability via machine learning, its role in
conditional trading strategies, and its behavior across realized volatility
regimes.

Empirical results confirm that the BVRP is positive on average, consistent
with the broader volatility risk premium literature, and exhibits negative
skewness and heavy tails reflecting episodic spikes in realized volatility.
OLS regressions with Newey--West standard errors show that the BVRP has
statistically significant predictive power for Bitcoin returns at 20- and
30-day horizons, though the explained variance remains modest, consistent with
the low signal-to-noise ratio of speculative asset markets.

Four supervised learning models---LASSO, Ridge, Random Forest, and Gradient
Boosting---are applied in a rolling out-of-sample design to forecast the level
of the BVRP. Tree-based models outperform linear benchmarks at specific
horizons, with realized volatility emerging as the dominant predictor,
accounting for approximately 69\% of total feature importance in the Random
Forest specification. BVRP quantile-based and regime-conditioned trading
strategies (Chapter 7) generate positive risk-adjusted returns, with
conditional rules filtering high-BVRP periods delivering the most consistent
performance.

A formal regime analysis (Chapter 9) classifies observations into low, medium,
and high realized volatility regimes using the 25th and 75th percentiles of
RV30D. The BVRP averages +13.25 percentage points in the low-volatility regime
and turns negative ($-1.50$ p.p.) in the high-volatility regime, reflecting
the tendency of realized volatility to exceed implied volatility during market
stress. Welch $t$-tests confirm that future returns differ significantly across
regimes at horizons of 10 days and beyond ($p < 0.05$), with the regime effect
operating additively rather than through a modulation of the BVRP's predictive
power.

Taken together, the results indicate that the BVRP functions primarily as a
volatility regime indicator. Its value for practitioners lies in conditioning
entry rules and risk filters rather than in pure directional forecasting.

\textbf{Keywords:} Bitcoin; implied volatility; realized volatility; volatility
risk premium; DVOL; Deribit; machine learning; random forest; gradient
boosting; volatility regimes; trading strategies.
